\documentclass[12pt,letterpaper]{article}
\usepackage[T1]{fontenc}
\usepackage{newtxtext,newtxmath}
\usepackage[letterpaper,margin=1in]{geometry}
\usepackage{amsmath,amssymb,microtype}
\usepackage[hidelinks]{hyperref}
\usepackage{xurl}
\title{Revision Log: Spectral Invariance under Grosswald--Schnitzer Prime Perturbations}
\author{David Joyner\\\texttt{wdjoyner@gmail.com}}
\date{Draft 2, November 28, 2026}
\begin{document}
\maketitle

\section{History}
\begin{itemize}
\item \emph{November 1, 2026.}  Research note with the same title,
  prepared by ChatGPT (OpenAI) for the author.  It is cited as CJ26 in
  Draft~2 of the companion paper
  \nolinkurl{prime_perturbations_revised.tex}.
\item \emph{November 28, 2026, unnumbered expanded revision.}  ChatGPT
  reformatted the paper for 12-point Times and letter paper.  It added an
  Introduction with a literature review and notation, expanded the
  technical glossary, and added the section on the November 2026 RH
  formalization with explicit limits.  The existing operator-theoretic
  claims were not newly certified.  (This entry replaces the former file
  \nolinkurl{revision_log_2026.md}, which covered both papers in this
  folder.)
\item \emph{November 28, 2026, Draft~2.}  Prepared by Claude (Anthropic)
  after a referee pass; details below.
\end{itemize}

\section{Draft 2: verification}
Claude checked the proofs of Lemma~3.1 (sublinear growth), Lemma~4.1
($Q_A$ is outer), Lemma~4.2 (ranges and the reproducing-kernel graph
core), Theorem~2.1 (equality of the closed joint kernels), and
Propositions~6.1, 6.2, 7.1 and~8.1, and found them correct.  The
isomorphism $U_k$ was previously cited to the earlier note without proof.
Claude verified that it is an isomorphism, and a short proof is now in
Section~5.  No mathematical error was found in the statements.  The
originality of the results has not been assessed.

\section{Draft 2: changes}
\begin{itemize}
\item \emph{RH consequences.}  The former Section~11 said that RH ``does
  not, by itself, prove spectral synthesis or a statement about the closed
  joint kernel''.  Theorem~2.1 of this paper is a spectral-synthesis
  theorem, so the new Corollary~11.1 states what follows.  Unconditionally,
  the closed joint kernel on $S_1$ vanishes if and only if RH holds; it is
  nonzero on $S_k$ for $k\ge2$.  With the certified RH, the closed joint
  kernel is $\{0\}$ on $S_1$ and is the closed span of the zero powers of
  all nontrivial zeros on $S_k$, $k\ge2$.  The abstract and Introduction
  were updated to match.
\item \emph{Graph section.}  The section ``Graph-theoretic interpretation:
  prime flows, boundaries, and transfer operators'' was a verbatim copy of
  Section~6 of the companion paper.  It used notation ($\sim_\delta$,
  $R_{e/f}$, $Q_e$) and results that appear only there.  It is replaced by
  a short Section~9.  That section states the boundary identity, gives
  Example~9.1, and cites the companion paper's classification.  The
  orphaned file \nolinkurl{gs_invariance_summation_operators_graph_appendix.tex},
  an older copy of the same material not input anywhere, was deleted.
\item \emph{Cross-references.}  Open question~1 of Section~10 now records
  that the companion paper answers it for endpoint shifts and for
  independent frequencies.  The citation of the ``earlier spectral-synthesis
  note'' now names it.  The companion paper is cited as PP26.
\item \emph{The order of the Sobolev norm.}  Remark~2.4 records that the
  corrected 1981 edition uses derivatives up to a fixed order $N$ rather
  than $k$, and that every argument goes through with $(\Delta+c)^N$ in
  place of $(\Delta+c)^k$.
\item \emph{Examples.}  The style guide asks for an example in every
  section.  Added Example~2.3 (moving only the first prime), Example~4.3
  (outer functions, bounded and unbounded), Example~5.1 (the first zeta
  zero lies in the joint kernel for $k=2$ but not in $S_1$), and
  Example~9.1 (telescoping along the prime ray).  New items were placed
  after existing numbered statements, so Theorem~2.1, Theorem~2.2,
  Lemmas~3.1--4.2 and Propositions~6.1--8.1 keep their numbers.
\item \emph{Remark after Proposition~7.1.}  The sentence ``the quotient
  can converge uniformly to $1$ on compact sets'' now says what varies: the
  size of the prescribed intervals.
\item \emph{Typesetting.}  Fixed \verb|$H^infty$|, \verb|$Re s$| and the
  graph-norm bars in the glossary.  Alphabetized the glossary, added the
  missing terms (dilation generator, Euler quotient, graph core,
  Grosswald--Schnitzer perturbation, $H^\infty$, Kronecker approximation,
  Laplace transform, reproducing kernel), and removed entries for terms
  no longer used.  Added a table of contents, the draft number, and
  section-based equation numbering.
\item \emph{Citations and attribution.}  The RH result is attributed as
  found with Anthropic's Claude, citing Rubinstein's page; the Lean FRO and
  Mathlib statement does not mention Claude.  The title of the Rubinstein
  citation was corrected, and an access date added.  The Acknowledgments
  now state each AI system's role.
\end{itemize}

\section{AI assistance}
ChatGPT (OpenAI) prepared the November~1 note and the November~28
expanded revision.  Claude (Anthropic) refereed the expanded revision and
prepared Draft~2.  AI review is not independent verification of the
mathematical research.
\end{document}
